\section{Boxes in the Store-Passing Interpreter}
\label{sec:storeboxes}

\subsection{The Box Clauses}

With the store in place, the four new clauses are short. \code{boxC} evaluates the initial contents, takes a fresh location, and adds a cell for it; \code{unboxC} evaluates the box and fetches its location:
\begin{lstlisting}
    [boxC (e)
          (type-case Result (eval-env env sto e)
            [res (v sto1)
                 (let ([l (new-loc)])
                   (res (boxV l) (override-store (cell l v) sto1)))])]
    [unboxC (e)
            (type-case Result (eval-env env sto e)
              [res (v1 sto1)
                   (let ([v2 (fetch (boxV-l v1) sto1)])
                     (res v2 sto1))])]
\end{lstlisting}
Note that \code{unboxC} fetches from \code{sto1}, the store \emph{after} evaluating \code{e}, in case evaluating the box expression itself wrote to it. \code{setboxC} evaluates the box and then the new contents, threading the store as \code{plusC} does, and adds a cell for the box's location on top of \code{sto2}:
\begin{lstlisting}
    [setboxC (e1 e2)
             (type-case Result (eval-env env sto e1)
               [res (v1 sto1)
                    (type-case Result (eval-env env sto1 e2)
                      [res (v2 sto2)
                           (res v1 (override-store (cell (boxV-l v1) v2) sto2))])])]
    [seqC (e1 e2)
          (type-case Result (eval-env env sto e1)
            [res (v1 sto1)
                 (eval-env env sto1 e2)])]
\end{lstlisting}
Two details:
\begin{itemize}
  \item \code{setboxC} returns \code{v1}, the box itself, not the new contents and not a void value (the language has none). This is a design choice; Racket's \code{set-box!} returns \code{(void)}.
  \item \code{seqC} discards \code{v1} but \emph{not} \code{sto1}: \code{e2} runs in the store that \code{e1} left behind. In Attempt~1 the \code{seqC} clause looked the same, but there was no store to pass on, so \code{e1} could only be evaluated and forgotten. Here sequencing finally has a purpose: \code{e1} is run for its effect on the store.
\end{itemize}

\subsection{Running the Examples}

Each in a fresh session, so that \code{new-loc} starts at 0:
\begin{lstlisting}[style=repl]
> (eval ex1)
- Result
(res (numV 4) '())
> (eval ex2)
- Result
(res (boxV 0) (list (cell 0 (numV 1)) (cell 0 (numV 0))))
\end{lstlisting}
\begin{lstlisting}[style=repl]
> (eval ex3)
- Result
(res (numV 0) (list (cell 0 (numV 0))))
> (eval ex4)
- Result
(res (numV 2) (list (cell 0 (numV 2)) (cell 0 (numV 1)) (cell 0 (numV 0))))
> (eval ex5)
- Result
(res (numV 12) (list (cell 0 (numV 2)) (cell 0 (numV 1))))
\end{lstlisting}
All five are now right. \code{ex1} allocates nothing, so its final store is empty. \code{ex2} is the result shown in lecture: the value is the box, location 0, and the store holds two cells for location 0, the write of 1 in front of the initial 0.

\subsection{Tracing \texttt{ex4}}

Write the store newest first, with \code{l:v} for \code{(cell l v)}. The outer \code{let} becomes an application of \code{(lambda b ...)} to \code{(box 0)}:
\begin{lstlisting}[style=repl]
(let b (box 0) (begin (begin (set-box! b 1)
                             (set-box! b (+ 1 (unbox b))))
                      (unbox b)))
  (lambda b ...)                          ---> closure, sto = []
  (box 0)       new location 0            sto = [0:0]
  body in env [b=(boxV 0)]
    (set-box! b 1)                        sto = [0:1, 0:0]
    (set-box! b (+ 1 (unbox b)))
      (unbox b): fetch 0 ---> 1
      (+ 1 1) ---> 2                      sto = [0:2, 0:1, 0:0]
    (unbox b): fetch 0 ---> 2
---> (numV 2) with sto = [0:2, 0:1, 0:0]
\end{lstlisting}
The environment never changes after the \code{let}: \code{b} is bound to \code{(boxV 0)} throughout. Every change happens in the store, and each \code{begin} hands the store from its first part to its second.

\subsection{Tracing \texttt{ex5}: Two Tables, Two Paths}

\begin{lstlisting}[style=repl]
(let a (box 1)
  (let f (lambda x (+ x (unbox a)))
    (begin (set-box! a 2) (f 10))))
  (box 1)       new location 0            sto = [0:1]
  env [a=(boxV 0)]
  (lambda x ...) ---> f = (closV [a=(boxV 0)] x (+ x (unbox a)))
  env [f=closure, a=(boxV 0)]
  (set-box! a 2)                          sto = [0:2, 0:1]
  (f 10)
    body env:   the closure's [a=(boxV 0)] + x=10
    body store: the current one,          sto = [0:2, 0:1]
    (unbox a): fetch 0 ---> 2
    (+ 10 2) ---> 12
---> (numV 12) with sto = [0:2, 0:1]
\end{lstlisting}
The closure saved the binding \code{a=(boxV 0)}: a location, not a number. The location has not changed since the closure was created, but the store has, and the call reads the current store (\Cref{fig:envstore}). This is how lexical scope and mutation coexist: the closure fixes \emph{which} box \code{a} is, and the store says what is in it \emph{now}.

\begin{figure}[htbp]
\centering
\begin{tikzpicture}[x=1cm, y=1cm]
  \node[cell, fill=tsk, minimum width=10mm] (cenv) at (0,0) {};
  \node[cell, fill=tsk, minimum width=7mm, right=0pt of cenv] (cx) {x};
  \node[cell, fill=tsk, minimum width=30mm, right=0pt of cx] (ce) {(+ x (unbox a))};
  \node[lbl, above=1mm of cx] {\code{f} = \code{closV}};
  \node[cell, fill=uspace, minimum width=22mm] (a) at (0,-1.6) {a=(boxV 0)};
  \fill (cenv.center) circle (0.6mm);
  \draw[ptr] (cenv.center) -- (a.north);
  \node[lbl, left=2mm of a, align=right, text=accent] {saved when the\\ \code{lambda} was evaluated};
  \node[cell, fill=kspace, minimum width=14mm] (x) at (0,-3.2) {x=10};
  \draw[ptr] (x) -- (a);
  \node[lbl, below=1mm of x, align=center] {environment of the body:\\ the closure's env + the parameter};
  \node[cell, fill=greybg, minimum width=12mm] (s2) at (5.6,-1.6) {0:2};
  \node[cell, fill=greybg, minimum width=12mm, right=5mm of s2] (s1) {0:1};
  \node[right=2mm of s1, lbl] (nil) {\code{[]}};
  \draw[ptr] (s2) -- (s1); \draw[ptr] (s1) -- (nil);
  \node[lbl, above=1mm of s2, xshift=8mm] {store \code{sto2}, passed in at the call};
  \draw[ptr, accent] (a.east) -- node[above, lbl, text=accent] {\code{fetch 0} stops here} (s2.west);
  \draw[decorate, decoration={brace, mirror, amplitude=4pt}] ($(s1.south west)+(0,-1mm)$) -- ($(nil.south east)+(0,-1mm)$) node[midway, below=2mm, lbl, align=center] {the store when\\ \code{f} was created};
\end{tikzpicture}
\caption{\code{ex5} while evaluating \code{(unbox a)} in \code{f}'s body. The environment (left) comes from the closure; the store (right) is the current one. The store at closure-creation time is still there as a suffix, but nothing reads it, since \code{fetch} stops at the newer cell.}
\label{fig:envstore}
\end{figure}

\begin{keypoint}[Where the environment and the store come from]
\begin{center}
\begin{tabular}{@{}p{0.3\linewidth}p{0.3\linewidth}p{0.3\linewidth}@{}}
\toprule
& environment & store \\
\midrule
to subexpressions & the same \code{env} to each & chained: \code{sto}, then \code{sto1}, \ldots \\
returned by \code{eval-env} & never & always, in the \code{res} \\
saved in a closure & yes & no \\
used for a function body & the closure's, extended & the current one, \code{sto2} \\
follows & the program text & the order of evaluation \\
\bottomrule
\end{tabular}
\end{center}
\end{keypoint}

\subsection{The Store Only Grows}

Nothing ever removes a cell. \code{set-box!} conses a new cell for an existing location, and the old one stays behind it, unreachable by \code{fetch}: after \code{ex4} the store holds three cells for location 0, of which only the first matters. Boxes that no part of the program can reach any more also keep their cells. A real implementation would overwrite a location in place and reclaim the memory of unreachable boxes with a garbage collector~\cite[Chapter~11]{PLAI}. For explaining what mutation \emph{means}, a store that only grows is simpler: every past state is still visible in the trace, and the store after any step is the store before it with some cells added on the front.
